\notechapter{N-20260405}{Computing Double Integrals: Fubini, Polar Coordinates, and Change of Variables}


\section{Iterated integrals}

If a plane region is vertically simple,
\[
  D=\{(x,y):a\le x\le b,\ \varphi_1(x)\le y\le \varphi_2(x)\},
\]
then Fubini's theorem gives
\[
  \iint_D f(x,y)\,dA=
  \int_a^b\int_{\varphi_1(x)}^{\varphi_2(x)}f(x,y)\,dy\,dx.
\]
If the region is horizontally simple, reverse the order.  The main work is not integrating; it is describing $D$ correctly.

\section{Changing the order}

Changing order means re-describing the same region with the other variable outermost.  A reliable method is: draw the projection interval, write lower and upper bounds, and split the region if a single formula cannot describe it.

\section{Polar coordinates}

For rotational or radial regions,
\[
  x=r\cos\theta,\qquad y=r\sin\theta,\qquad dA=r\,dr\,d\theta.
\]
The factor $r$ is the Jacobian.  It is the area scaling of the map $(r,\theta)\mapsto(x,y)$.

\section{General change of variables}

For a $C^1$ one-to-one transformation
\[
  x=x(u,v),\qquad y=y(u,v),
\]
the formula is
\[
  \iint_D f(x,y)\,dx\,dy=
  \iint_{D'} f(x(u,v),y(u,v))
  \left|\frac{\partial(x,y)}{\partial(u,v)}\right|\,du\,dv.
\]
The determinant is the local area multiplier.  This is the conceptual reason every coordinate transformation carries a Jacobian.

\section{The conceptual payoff}

Double-integral computation is a problem of choosing coordinates that make the region and integrand simple at the same time.  Fubini slices the region; polar coordinates exploit rotation; general change of variables is local linearization of area.



\paragraph{Expanded synthesis.}
For this chapter, the elementary material should be read as a study of what remains invariant when coordinates change. The earlier sections (`Iterated integrals`, `Changing the order`, `Polar coordinates`, `General change of variables`, `Higher viewpoint`) compute with arrays, bases, pivots, determinants, or eigenvectors; the deeper object is a linear map together with the spaces it acts on. A matrix is a representation of that map after bases have been chosen. This is why formulas such as
\[
  A' = P^{-1}AP,\qquad \widetilde A = T^{-1}AS
\]
are not cosmetic: they distinguish an intrinsic morphism from a coordinate description.

The structural hierarchy is as follows. Kernels record what a map forgets, images record what it reaches, cokernels record obstructions to solvability, and quotients record information after an equivalence has been imposed. Determinants act on the top exterior power and therefore measure oriented volume; traces are invariant under cyclic rearrangement and become characters in representation theory; adjoints depend on an inner product and lead to self-adjoint spectral theory.

A useful way to return to the computations is to ask, for every formula, which invariant it is measuring. Row reduction measures rank and kernel; diagonalization measures decomposition into invariant subspaces; Gram--Schmidt measures orthogonal projection; positive definiteness measures ellipsoid geometry; tensor notation measures how components transform. The elementary methods are therefore the finite-dimensional laboratory in which duality, quotienting, functoriality, and spectral decomposition can be seen clearly.


\section{Expanded order-changing method}

To change an integral order, never start by manipulating symbols.  Start from the region.  For example, if
\[
  0\le x\le1,\qquad x^2\le y\le x,
\]
then the projection on the $y$-axis is $0\le y\le1$, and the horizontal slice satisfies
\[
  y\le x\le \sqrt y.
\]
Thus
\[
  \int_0^1\int_{x^2}^{x}f(x,y)\,dy\,dx
  =
  \int_0^1\int_y^{\sqrt y}f(x,y)\,dx\,dy.
\]
The method is always geometric: inequalities must be solved as a description of the same set.

\section{Jacobian as local area determinant}

For $T(u,v)=(x(u,v),y(u,v))$, the derivative matrix
\[
  DT=\begin{pmatrix}x_u&x_v\\y_u&y_v\end{pmatrix}
\]
sends a tiny $du\,dv$ rectangle to an approximate parallelogram.  Its area is multiplied by $|\det DT|$.  This is why the absolute value of the Jacobian appears in change of variables.

\section{Jacobian as determinant of a pullback}

For a coordinate change
\[
  x=x(u,v),\qquad y=y(u,v),
\]
the differential forms transform as
\[
  dx=x_u\,du+x_v\,dv,\qquad dy=y_u\,du+y_v\,dv.
\]
Taking the wedge product gives
\[
  dx\wedge dy=(x_uy_v-x_vy_u)\,du\wedge dv.
\]
Thus
\[
  dx\wedge dy=\det\frac{\partial(x,y)}{\partial(u,v)}\,du\wedge dv.
\]
The Jacobian is the coefficient by which the coordinate change pulls back the area form.

This is the differential-form version of the elementary change-of-variables formula.

\section{Coarea formula preview}

Changing variables is one way to decompose integration.  The coarea formula is a more advanced generalization.  For a smooth map $F:\mathbb R^n\to\mathbb R^k$, it relates an integral over the domain to integrals over level sets of $F$:
\[
  \int_{\mathbb R^n} g(x)J_F(x)\,dx
  =\int_{\mathbb R^k}\left(\int_{F^{-1}(y)} g\,d\mathcal H^{n-k}\right)dy.
\]
For $k=1$, this decomposes volume by level surfaces.

The elementary method of slicing a region vertically or horizontally is a simple coarea idea: integrate first over fibers, then over the parameter indexing the fibers.

\section{Orientation and absolute value}

In ordinary area integrals, the absolute value of the Jacobian appears:
\[
  dA=|J|\,du\,dv.
\]
For oriented forms, the sign is retained:
\[
  dx\wedge dy=J\,du\wedge dv.
\]
This distinction explains why flux integrals and Green/Stokes formulas care about orientation, while scalar area integrals do not.

\section{Returning to order changes}

Changing integration order is not a syntactic operation on bounds; it is a choice of projection and fibers.  General change of variables replaces rectangular fibers by smooth coordinate fibers.  The original exercises teach the finite-dimensional geometry needed for the more general integration theorems.

% Second-pass deepening for waves 2-4: wave 4 part 2
\section{Change of variables as local linearization}

At a point $u$, the differentiable map $T$ is approximated by its derivative
\[
  T(u+h)=T(u)+DT(u)h+o(|h|).
\]
Small squares in parameter space become approximate parallelograms in physical space.  The determinant $\det DT(u)$ measures their oriented area scaling.

This is why the Jacobian appears locally and is then integrated globally.  The change-of-variables theorem is the accumulation of infinitely many local linear approximations.

\section{Singular coordinate systems}

Polar coordinates fail to be one-to-one at the origin and are periodic in $\theta$.  Spherical coordinates have singularities at the poles.  These singularities usually do not harm integrals because they occur on sets of measure zero, but they matter for topology and vector fields.

This explains why coordinate formulas must be used with awareness of domains.  A coordinate system may be excellent for integration but poor for global topology.
